\documentclass{article}

\usepackage{amsmath, amsthm, amssymb, amsfonts}
\usepackage{thmtools}
\usepackage{graphicx}
\usepackage{setspace}
\usepackage{geometry}
\usepackage{float}
\usepackage{hyperref}
\usepackage[utf8]{inputenc}
\usepackage[english]{babel}
\usepackage{framed}
\usepackage[dvipsnames]{xcolor}
\usepackage{tcolorbox}
\usepackage{physics}
\usepackage[shortlabels]{enumitem}

\DeclareMathOperator{\grd}{grad}
\DeclareMathOperator{\crl}{curl}
\DeclareMathOperator{\dvr}{div}
\DeclareMathOperator{\length}{length}
\DeclareMathOperator{\vspan}{span}
\DeclareMathOperator{\rng}{range}
\DeclareMathOperator{\nul}{null}
\colorlet{LightGray}{White!90!Periwinkle}
\colorlet{LightOrange}{Orange!15}
\colorlet{LightGreen}{Green!15}

\newcommand{\HRule}[1]{\rule{\linewidth}{#1}}

\declaretheoremstyle[name=Theorem,]{thmsty}
\declaretheorem[style=thmsty,numberwithin=section]{theorem}
\tcolorboxenvironment{theorem}{colback=LightGray}

\declaretheoremstyle[name=Proposition,]{prosty}
\declaretheorem[style=prosty,numberlike=theorem]{proposition}
\tcolorboxenvironment{proposition}{colback=LightOrange}

\declaretheoremstyle[name=Principle,]{prcpsty}
\declaretheorem[style=prcpsty,numberlike=theorem]{principle}
\tcolorboxenvironment{principle}{colback=LightGreen}

\setstretch{1.2}
\geometry{
    textheight=9in,
    textwidth=6in,
    top=0.5in,
    headheight=12pt,
    headsep=25pt,
    footskip=30pt
}

% ------------------------------------------------------------------------------

\begin{document}


% ------------------------'

\textbf{Final Exam}

\textbf{Excercise 1.} Find the change of coordinates matrix from the basis $1$, $1+t$ to the basis $1-t$, $2t$ in $\mathbb{P}_1$.

\textbf{Excercise 2.} Let $Perm(n)$ be the set of permutations of $\{1,2,\ldots,n\}$. Let $\tau \in Perm(n)$. Show that the map from $Perm(n)$ to $Perm(n)$ defined by $\sigma \mapsto \sigma \tau$ is a bijection (one-to-one and onto). In other words, composition with a fixed permutation $\tau$ takes the list of all permutations and rearranges it, but does not duplicate or remove any (it permutes the permutations).

\textbf{Excercise 3.} Show that the determinant of a matrix is the product of its eigenvalues (counting multiplicities). 
\textbf{Hint:} First show that the characteristic polynomial $\det(A-\lambda I) = (\lambda_1 - \lambda)(\lambda_2 - \lambda)\cdots(\lambda_n - \lambda)$ (note the $\lambda_i$'s occur before $\lambda$), where $\lambda_1, \ldots, \lambda_n$ are the eigenvalues of $A$ (with multiplicity). Then plug in $\lambda = 0$. Note that in general, a polynomial of degree $n$ can be written as $p(\lambda) = c(\lambda - r_1)(\lambda - r_2)\cdots(\lambda - r_n)$ (note the $\lambda$ comes before each $\lambda_i$), so the first step is to show that this $c$ is $(-1)^n$. If you use the permutation definition of the determinant, there is exactly one term that is of degree $n$ in $\lambda$, and it comes from the identity permutation.

\textbf{Excercise 4.} Suppose the entries of $A$ and $A^{-1}$ are integers. What are the possible values of $\det(A)$? \textbf{Hint:} Use the fact that $\det(A) \cdot \det(A^{-1}) = 1$.


\end{document}